% Biais inductifs des méthodes de flot optique : tableau simple, police Times
\documentclass[border=10pt,varwidth=900pt]{standalone}
\usepackage[T1]{fontenc}
\usepackage[utf8]{inputenc}
\usepackage{mathptmx}
\usepackage{array}
\usepackage{hhline}
\usepackage[table]{xcolor}
\renewcommand{\arraystretch}{1.45}
\setlength{\arrayrulewidth}{0.6pt}
\definecolor{rougep}{HTML}{FAD4D4}
\definecolor{vertp}{HTML}{D6EFD0}
\newcommand{\oui}{\cellcolor{rougep}oui}
\newcommand{\non}{\cellcolor{vertp}non}
\newcolumntype{C}{>{\centering\arraybackslash}m{78pt}}

\begin{document}
\fontsize{14}{17}\selectfont
\noindent
\begin{tabular}{|>{\raggedright\arraybackslash}m{115pt}|C|C|C|C|C|C|>{\centering\arraybackslash}m{100pt}|}
\hhline{|-|-|-|-|-|-|-|-|}
\textbf{Méthode} & \textbf{Corrélation} & \textbf{Warping} & \textbf{Pyramide} & \textbf{Itérations} & \textbf{Upsampling convexe} & \textbf{Tiling} & \textbf{Tête} \\
\hhline{|-|-|-|-|-|-|-|-|}
PWC-Net & \oui & \oui & \oui & \non & \non & \non & --- \\
\hhline{|-|-|-|-|-|-|-|-|}
RAFT & \oui & \non & \non & \oui & \oui & \non & --- \\
\hhline{|-|-|-|-|-|-|-|-|}
MemFlow & \oui & \non & \non & \oui & \oui & \non & --- \\
\hhline{|-|-|-|-|-|-|-|-|}
FlowFormer & \oui & \non & \non & \oui & \oui & \oui & --- \\
\hhline{|-|-|-|-|-|-|-|-|}
WAFT / GeoViT & \non & \oui & \non & \oui & \oui & non / oui & --- \\
\hhline{|-|-|-|-|-|-|-|-|}
CroCo-Flow & \non & \non & \non & \non & \non & \oui & DPT \\
\hhline{|-|-|-|-|-|-|-|-|}
\textbf{FreeFlow} & \non & \non & \non & \non & \non & \non & 3 convolutions \\
\hhline{|-|-|-|-|-|-|-|-|}
\end{tabular}

\vspace{8pt}
\noindent\colorbox{rougep}{\strut\ oui\ } la méthode utilise ce module fait à la main \qquad
\colorbox{vertp}{\strut\ non\ } elle ne l'utilise pas
\end{document}
